\section{Related Work}
\label{app:related_work}
\subsection{Cognitive Science Based Agent Memory}
Cognitive science has long studied the mechanisms of human memory, providing valuable insights into how memory can be leveraged in agents. The Atkinson–Shiffrin memory model defines long-term memory as information that can be retained and recalled over extended periods of time \citep{Atkinson1968HumanMemory}. Long-term memory is further categorized into three major types \citep{Tulving1972EpisodicSemantic, Squire2004MemorySystems}: semantic memory, which stores concepts and factual knowledge; procedural memory, which is responsible for knowing how to perform actions and skills; and episodic memory, which represents raw records of personal experiences and events. For intelligent agents, most interactions with the environment are episodic memory, such as raw contexts and working trajectories. From these episodic experiences, agents can extract more reusable representations. For example, user preferences can be distilled as semantic memory \citep{li2025memosmemoryosai}, while web navigation skills can be abstracted as procedural memory \citep{wang2024agentworkflowmemory}. In this sense, semantic and procedural memories are not isolated components, but rather structured abstractions derived from episodic memory. Knowledge can be viewed as a further abstraction over memory. Specifically, propositional knowledge corresponds to highly structured and verifiable abstractions of semantic memory, while prescriptive knowledge represents generalized and reusable forms of procedural memory that describe how to accomplish goal-oriented tasks \citep{10.1257/aer.103.3.545}. \par
Inspired by these cognitive theories, many studies have explored memory mechanisms for intelligent agents. Episodic and semantic memory have been introduced to handle long-context information, such as long documents in document understanding agents \citep{lee2024humaninspiredreadingagentgist} and long-term user–agent interactions in conversational agents \citep{li2025memosmemoryosai}. Additionally, some approaches construct knowledge graphs containing both episodic and semantic subgraphs to better organize and manage agent memory \citep{anokhin2024arigraphlearningknowledgegraph,rasmussen2025zeptemporalknowledgegraph}. Procedural memory has also been incorporated into decision-making agents, particularly web agents, where accumulated experiences are leveraged to improve future strategies \citep{zhao2024expelllmagentsexperiential}. This line of research moves toward self-evolving agents that continuously refine their decision-making policies based on incoming memory \citep{wang2024agentworkflowmemory, DBLP:journals/corr/abs-2508-06433, ouyang2025reasoningbankscalingagentselfevolving}. \par
However, most existing methods are designed for specific tasks and are unable to effectively support all memory types simultaneously. In contrast, \plugmem standardizes heterogeneous episodic memories across diverse tasks and stores extracted knowledge in a propositional–prescriptive dual knowledge graph, thereby overcoming the limitations of task generalizability present in prior approaches.

\subsection{Memory Module Designs}

Across agentic systems, external memory is most commonly realized through retrieval. Agents store past interactions, observations, or documents outside the model context and retrieve relevant information at inference time to condition decision making. This retrieval first paradigm forms the basic prototype of external memory, and nearly all subsequent memory designs can be viewed as extensions built on top of it.

Early task agnostic memory systems adopt flat retrieval over unstructured episodic memory. Vanilla retrieval and vanilla RAG store raw interaction histories or documents as text and retrieve relevant segments based on similarity to the current query~\citep{lewis2021retrievalaugmentedgenerationknowledgeintensivenlp,Wang_2024}. These methods provide a simple and broadly applicable memory interface, but the retrieved content remains tightly bound to specific episodes. As a result, retrieved memories often contain redundant details and exhibit low information density when reused across different decision contexts.

To improve retrieval effectiveness, later work introduces structure on top of episodic memory. Graph based and hierarchical retrieval systems organize stored content to support multi hop retrieval, aggregation, or global reasoning~\citep{edge2025localglobalgraphrag,sarthi2024raptorrecursiveabstractiveprocessing}. Related approaches further refine retrieval by altering the internal organization of memory graphs, such as ranking paths or linking propositions~\citep{gutiérrez2025ragmemorynonparametriccontinual,Wang_2025}. While these methods differ in structure and traversal strategy, they largely preserve the same retrieval oriented abstraction. Memory remains organized around episodic traces, entities, or text spans, and structural mechanisms primarily serve to improve access to stored experiences rather than to transform their content.

In parallel, task-specific memory systems explore more aggressive forms of memory processing. Temporal knowledge graphs for conversational agents, workflow memories for web navigation, and reasoning oriented memory banks explicitly extract higher level representations from experience~\citep{gutiérrez2025ragmemorynonparametriccontinual,wang2024agentworkflowmemory,ouyang2025reasoningbankscalingagentselfevolving}. These systems demonstrate that abstracting experience into reusable representations can substantially improve performance within a fixed task setting. However, the abstraction process is typically guided by task-specific assumptions, which constrains reuse of the memory module across heterogeneous agents and environments.

Viewed together, prior work reveals a key distinction in how memory is treated. Some systems focus on improving retrieval over episodic experience, while others implicitly or explicitly transform experience into more abstract forms. Only the latter enables memory to generalize across episodes and tasks, since episodic memory is inherently tied to a particular agent and interaction trajectory, whereas semantic and procedural knowledge abstract away task-specific execution details. This observation motivates memory to knowledge transformation as a critical operation for reusable agent memory.

\plugmem builds on this evolution by making memory-to-knowledge transformation an explicit and central component of the memory module. Instead of treating episodic traces as the primary unit of retrieval, \plugmem organizes memory around propositional and prescriptive knowledge units corresponding to semantic and procedural memory. These knowledge units serve as the basis for memory organization, retrieval, and reasoning, while episodic memory is retained as verifiable source evidence. By elevating knowledge to the unit of memory manipulation, \plugmem enables memory to be reused across agents and tasks without task-specific redesign, while remaining compatible with the retrieval based external memory paradigm.


\subsection{Memory Benchmarks and Evaluation Scope}

A broad range of benchmarks have been used to study memory related behaviors in language models and agents, spanning diverse task settings and evaluation objectives. One prominent line of work focuses on long context utilization, evaluating whether models can access and reason over large inputs within a single inference window, such as long document understanding, narrative comprehension, or needle in a haystack style retrieval~\citep{liu2023lostmiddlelanguagemodels,bai2024longbenchbilingualmultitaskbenchmark,zhang2024inftybenchextendinglongcontext}. These benchmarks assess context usage and attention behavior, but do not model memory accumulation, update, or reuse across time.
A separate line of work focuses on long-horizon conversational agents and user-centric memory, evaluating whether systems can accumulate, maintain, and reuse information such as preferences, personal attributes, and dialogue context over extended interactions~\citep{wu2024longmemeval,maharana2024evaluatinglongtermconversationalmemory}. These benchmarks explicitly stress episodic accumulation and robustness to redundancy and noise, and are commonly used to study memory abstraction in interactive settings.
Another class of benchmarks evaluates retrieval-augmented question answering over static knowledge sources, including factoid and multi-hop QA tasks, where memory is treated as a fixed external repository and the primary challenge lies in retrieval and evidence composition rather than memory evolution~\citep{petroni2021kiltbenchmarkknowledgeintensive,kwiatkowski-etal-2019-natural,yang2018hotpotqadatasetdiverseexplainable,trivedi2022musiquemultihopquestionssinglehop}. While effective for studying semantic retrieval and composition, these benchmarks assume a static memory base and do not capture how memory is constructed or refined through interaction.
In agentic settings, additional benchmarks evaluate end-to-end task performance in interactive environments, including web navigation, tool use, and embodied reasoning, where memory interacts with planning, perception, and action execution~\citep{zhou2024webarenarealisticwebenvironment,trivedi2024appworldcontrollableworldapps,furuta2024multimodalwebnavigationinstructionfinetuned,shridhar2021alfworldaligningtextembodied,liu2025agentbenchevaluatingllmsagents}. In such benchmarks, memory is often entangled with other agent capabilities, making it difficult to isolate memory-specific contributions.

While these benchmarks differ widely in surface form, they each probe only a limited slice of memory behavior. Collectively, prior work reveals a fragmented evaluation landscape, where no single benchmark provides comprehensive coverage of how memory is accumulated, abstracted, organized, and reused in agent decision-making.

Within this landscape, memory benchmarks can be distinguished by how memory is constructed, organized, and reused. Some benchmarks emphasize long term accumulation of episodic experience, requiring agents to extract stable information from extended interaction histories. Others focus on structured access to semantic knowledge, stressing multi-hop retrieval and organization over memory. A third class emphasizes procedural reuse, where agents must generalize skills or workflows acquired from prior experience to new task instances. No single benchmark captures all of these aspects simultaneously, motivating the need for a complementary evaluation suite.

Based on this perspective, we intentionally select three benchmarks that probe distinct and complementary memory behaviors. LongMemEval targets long-horizon conversational settings where memory is accumulated over time and relevant information must be abstracted from noisy episodic interactions~\citep{wu2024longmemeval}. HotpotQA evaluates semantic memory organization and multi-hop retrieval over a large static knowledge source, serving as a canonical benchmark for structured memory access and composition~\citep{yang2018hotpotqadatasetdiverseexplainable}. WebArena evaluates procedural memory in interactive environments, where agents must reuse and adapt previously acquired strategies across multiple task instances derived from shared intent templates~\citep{zhou2024webarenarealisticwebenvironment}.

Beyond their complementary coverage of episodic accumulation, semantic organization, and procedural reuse, the selected benchmarks are representative within their respective evaluation settings. LongMemEval, HotpotQA, and WebArena have each been widely adopted in prior work, with a broad range of memory and agent baselines already evaluated on these benchmarks. This established usage allows existing systems to be reliably reproduced and compared under consistent protocols, reducing evaluation bias introduced by task-specific tuning. As a result, evaluating memory modules on these benchmarks enables fairer and more stable comparison across heterogeneous memory designs, while supporting conclusions that generalize beyond any single task or interaction pattern.
